\documentclass[10pt,a4paper]{amsart}
\usepackage[margin=19mm]{geometry}
\usepackage{amsmath,amssymb,mathtools}
\usepackage[scr=rsfs,cal=euler]{mathalfa}
\usepackage{xfrac}
\usepackage[unicode,hidelinks,bookmarks=false]{hyperref}
\newcommand{\ie}{i.e.\ }
\newcommand{\cf}{cf.\ }
\newcommand{\resp}{respectively\ }
\newcommand{\Subsection}{\subsection}
\providecommand{\nobreakdash}{\nobreak\hbox{-}}
\setlength{\emergencystretch}{2em}
\multlinegap0pt
\usepackage{bm}
\usepackage{bbm}
\usepackage{dsfont}
\usepackage{xspace}
\usepackage{cancel}
\usepackage{mathtools}
\usepackage{amsthm}
\makeatletter
\@ifundefined{theo}{\newtheorem{theo}{Theo}}{}
\@ifundefined{lemm}{\newtheorem{lemm}[theo]{Lemma}}{}
\makeatother
\title[Untriangular factorization of holomorpic symplectic matrices]{Untriangular factorization of holomorpic symplectic matrices}
\author{Gaofeng Huang, Frank Kutzschebauch, Phan Quoc Bao Tran}
\date{Source: arXiv:2507.18963v1 (CC BY 4.0)}
\begin{document}
\maketitle

\noindent\emph{Source and licence.} Untriangular factorization of holomorpic symplectic matrices. Version arXiv:2507.18963v1 (CC BY 4.0), distributed under the Creative Commons Attribution 4.0 International licence, as recorded in the arXiv licence metadata for this exact version and shown on the abstract page https://arxiv.org/abs/2507.18963v1. Full rights notice: licenses/arxiv-2507.18963-CC-BY-4.0.txt.\par\smallskip

\noindent\textit{Editorial scope.} Counted unit: the Lemm whose source label is \texttt{red-2} in the article (source0.tex lines 1581--1583, source label \texttt{red-2}), delivered together with the article's own complete original proof of it (source0.tex lines 1584--1660, one \texttt{proof} environment of 77 lines). No mathematics is written, repaired or re-derived: statement and proof are the article's own text line for line. Required context retained uncounted: equa-2 (lines 1494--1499); red-1 (lines 1522--1524). Every definition, display and label of those lines is retained unchanged. Omitted from this package and not redistributed: 1--1493, 1500--1521, 1525--1580, 1661--2334. Nothing inside a retained excerpt is omitted or altered: every retained body is delivered exactly as the article's lines are written, so no omission receipt of any kind accompanies this package. Citations: the retained excerpts carry no citation command at all, so no bibliography entry of the article is transcribed and no reference is redistributed. Reference adaptation: none was needed and none was made; every reference inside the delivered unit resolves inside it, so no unresolved reference or citation is produced. Printed slips kept literal: no printed word, symbol, hypothesis or number of the counted statement or of its proof was altered. Layout and class: the article's own class is \texttt{amsart} with packages including comment, amsmath, amsfonts, euscript, amscd, amsthm, amssymb, upref, graphics, mathrsfs, yhmath, xy, hyperref, amsrefs; the wrapper is the lane's pinned \texttt{amsart} editorial wrapper at 10pt on A4, which restates the theorem-like environments the retained text needs and adds the article's own macro definitions that the retained text needs (none), with extra wrapper packages (bm, bbm, dsfont, xspace, cancel, mathtools). The delivered numbering is the unit's own. Assets: no retained span contains a figure environment, a graphics-inclusion command, a PDF-page inclusion, a table or a TikZ picture; no image file is copied into this package, the package declares no asset dependency and no image file is redistributed with it. Licence: version arXiv:2507.18963v1 of this article is distributed under the Creative Commons Attribution 4.0 International licence, as recorded in the arXiv licence metadata for this exact version and shown on the abstract page https://arxiv.org/abs/2507.18963v1. A full rights notice is delivered at licenses/arxiv-2507.18963-CC-BY-4.0.txt. Characters: every retained excerpt is pure ASCII, so the delivered engine, XeLaTeX with its default Latin Modern text font, reports no missing character.
\begin{equation}\label{equa-2}
    \begin{cases}
    P^K_f A_K = P^{K-1}_{f} \\
      P^K_s A^{-T}_K = P^{K-1}_s + P^{K-1}_f Z_K 
\end{cases}
\end{equation}

\begin{lemm}\label{red-1}
Let $(a,b)\in G_i \subset \mathbb{C}^{2n}\setminus \{0\}, i = 1, \dots, n$. Then the fiber $(P^K)^{-1}(a,b)$ given by the set of $2n$ equations $P^K(A,Z) = (a, b)$ can be described by one equation, namely $P^K_i = a_i$.
\end{lemm}

\begin{lemm}\label{red-2}
     Let $(0,b)\in N_i \subset \mathbb{C}^{2n}\setminus \{0\}, i = 0, \dots, n-1$. Then the fiber $(\Phi_K|_{U_K})^{-1}(0,b)$ given by the set of $2n$ equations $P^K(A,Z)=(0,b)$ can be described by one equation, namely $P^K_{2n-i}=b_{n-i}$.
\end{lemm}

\begin{proof}%[Proof of Lemma \ref{red-2}:]
Recall the inductive formula
\begin{align*}
    P^K = P^{K-1} M^+(A_K,Z_K) = P^{K-1} \left[\begin{array}{ccc} I & Z_K \\ 0 & I\end{array}\right] \left[\begin{array}{ccc} A^{-1}_K & 0 \\ 0 & A^T_K\end{array}\right].
\end{align*}
We consider the fiber of $(0,b)$ in $N_i$. Then (\ref{equa-2}) becomes
\begin{align*}
    \begin{cases}
    0 = P^{K-1}_{f} \\
    b A^{-T}_K = P^{K-1}_s  
    \end{cases}
\end{align*}
which is equivalent to
\begin{align} \label{P_f=0}
    \begin{cases}
    0 = P^{K-1}_{f}  \\
    b  = P^{K-1}_s  A^{T}_K
    \end{cases}
\end{align}
Writing out the second equation, we have
\[
    ( b_1, \dots, b_{n-i}, 0, \dots, 0) = P^{K-1}_s 
    \begin{bmatrix}
     1 & 0 & \cdots &0 \\  a_{12,K} & 1 & \ddots & 0 \\
    \vdots & \ddots & \ddots & \vdots \\
    a_{1n,K} & \cdots & a_{n-1,n,K} & 1
    \end{bmatrix}
\]
The last $n+i+1$ components yield
\begin{align} \label{n+i+1}
    P^{K-1}_{2n-i}=b_{n-i}, \, P^{K-1}_{2n-i+1}=0, \, \dots, P^{K-1}_{2n-1} = 0, \, P^{K-1}_{2n} = 0.
\end{align}
Since $b_{n-i}$ is nonzero, we can use the first $n-i-1$ equations to express the variables $a_{k,n-i,K}, k = 1, \dots, n-i-1$ in $A_K$ as in Step 1 of the proof of Lemma \ref{red-1}.

We write $P^{K-1} = P^{K-2} M^-(A_{K-1}, Z_{K-1})$
\begin{center} \label{P^K-1}
    $P^{K-1} = P^{K-2} \left[\begin{array}{ccc} I_n & 0 \\ Z_{K-1} & I_n\end{array}\right] \left[\begin{array}{ccc} A^{-T}_{K-1} & 0 \\ 0 & A_{K-1}\end{array}\right] = P^{K-2} \left[\begin{array}{ccc} A^{-T}_{K-1} & 0 \\ 0 & A_{K-1}\end{array}\right] \left[\begin{array}{ccc} I_n & 0 \\ \tilde{Z}_{K-1} & I_n\end{array}\right]$,
\end{center}
where $\tilde{Z}_{K-1}:= A_{K-1}^{-1} Z_{K-1} A_{K-1}^{-T}$ is also symmetric. 


\begin{center}
    $P^{K-1} \left[\begin{array}{ccc} I_n & 0 \\ -\tilde{Z}_{K-1} & I_n\end{array}\right] = P^{K-2} \left[\begin{array}{ccc} A^{-T}_{K-1} & 0 \\ 0 & A_{K-1}\end{array}\right]$
\end{center}
equivalently
\begin{equation}
\begin{cases}
    P^{K-1}_f - P^{K-1}_s \tilde{Z}_{K-1} &= P^{K-2}_f A^{-T}_{K-1} \\
    P^{K-1}_{s} A^{-1}_{K-1} &= P^{K-2}_{s} 
\end{cases}  
\end{equation}
Therefore, the first set of equations from \eqref{P_f=0} and the equations \eqref{n+i+1} are equivalent to 
\begin{align*}
    \begin{cases}
        - P^{K-1}_s \tilde{Z}_{K-1} &= P^{K-2}_f A^{-T}_{K-1} \\
        P^{K-1}_{2n-i} &= b_{n-i} \\
        \sum^{j-1}_{k=1}d_{kj}P^{K-1}_{n+k} & = P^{K-2}_{n+j}, \quad j = n-i+1, \cdots, n  
    \end{cases}
\end{align*}
where we choose 
\begin{center}
     $A^{-1}_{K-1}  = \left[\begin{array}{cccc} 1 & d_{12} & \cdots & d_{1n} \\  0 & 1 & \ddots & \vdots \\
    \vdots & \ddots & \ddots & d_{n-1,n} \\
    0 & \cdots & 0 & 1\end{array}\right]$.
\end{center}
A closer inspection on the last set of equations, together with \eqref{n+i+1}, gives
\begin{align}
    d_{n-i,j} b_{n-i} + \sum^{n-i-1}_{k=1}d_{kj}P^{K-1}_{n+k} = P^{K-2}_{n+j}  \quad \text{ for } j = n-i+1, \dots, n. 
\end{align}
Since $b_{n-i}$ is nonzero, we can use these equations to express $d_{n-i,j}$
\[
     d_{n-i,j} =  (P^{K-2}_{n+j} - \sum^{n-i-1}_{k=1}d_{kj}P^{K-1}_{n+k})/b_{n-i}  \quad  \text{ for } j = n-i+1, \dots, n. 
\]
Since $P^{K-1}_s \neq 0$, the first set of equations $- P^{K-1}_s \tilde{Z}_{K-1} = P^{K-2}_f A^{-T}_{K-1}$ can be used to express $n$ variables in $\tilde{Z}_{K-1}$ as in Step 2 of proof of Lemma \ref{red-1}. 

The introduction of $\tilde{Z}_{K-1}$ and $d_{ij}, i>j$ is another choice of coordinates on the set of unitriangular $\tilde{\Omega}$-symplectic matrices. It allows to reduce the set of $2n$ equations in a simple way to one single equation. 
\end{proof}

\end{document}
